\documentstyle[12pt]{article}
\begin{document}

\begin{center}
{\large\bf Analysis of Bifurcation and Stability of Periodic Systems\\
 Using Truncated Point Mappings\\}
\vspace{0.25in}
{\sc Ramesh S. Guttalu and Henryk Flashner}\\
Department of Mechanical Engineering \\
University of Southern California\\
Los Angeles, CA 90089-1453
\end{center}
 \ \vspace{0.25in}
\begin{center}
{\large\bf ABSTRACT \\[2ex]}
\end{center}

An approach to analyze stability and bifurcation of nonlinear
systems with periodically varying coefficients is presented. The method is
based on a {\em point mapping} (period to period mapping) representation of
system's dynamics. An algorithm to obtain an analytical expression for the
point mapping and its dependence on system's parameters is employed. The
algorithm is devised to derive the coefficients of a multinomial expansion
of the point mapping up to an arbitrary order in terms of the state
variables and of the parameters. Analytical stability and bifurcation
condition are then formulated and expressed as functional relations between
the various parameters. Algebraic equations for determining bifurcated periodic
solutions of different periods are derived in analytical form, providing an
explicit dependence of these solutions on system's parameters. To
demonstrate the capabilities of the method, it is applied for investigating
the parametric stability and bifurcation of a nonlinear Mathieu's equation.
The results obtained by the proposed approach are compared to those obtained
by direct integration that are considered to be the ``exact solution''. It
is shown that, unlike perturbation analysis, the proposed method provides
very accurate solution even for {\em large values} of the parameters.
Moreover, the proposed framework can be easily applied to
multiple degree-of-freedom systems with multiple parameters.  This feature is
an important advantage of the proposed approach, since most of the
existing analysis methods apply mainly to single degree-of-freedom systems.
Extension of these methods to higher
dimensions is difficult and computationally cumbersome.

\end{document}
